Let \(m>1\) be an integer prime to the characteristic of \(k\).
By \SGARefLink{sga3:viii:corollary:6:7}{Exposé VIII, 6.7}, for every \(h>0\),
\(\operatorname{Centr}_{G_0}({}_{m^h}T_0)\) is represented by a
closed subscheme of \(G_0\).  The subgroup schemes
\({}_{m^h}T_0\) are schematically dense in \(T_0\)
\textup{(cf.\ \SGARefLink{sga3:ix:remark:4:10}{Exposé IX, 4.10})}, and \(G_0\) is noetherian.
Consequently there is an \(h\) such that
\[
  \operatorname{Centr}_{G_0}({}_{m^h}T_0)
  =
  \operatorname{Centr}_{G_0}(T_0)
  =
  T_0.
\]
Put \(n=m^h\).  Since \(n\) is invertible on \(S\),
\({}_n\mathbb G_{m,S}\) is isomorphic to
\((\mathbb Z/n\mathbb Z)_S\), so \(f_0\) defines a group
monomorphism
\[
  u_0\colon(\mathbb Z/n\mathbb Z)_k^{\,r}\longrightarrow G_0
\]
such that \(\operatorname{Centr}_{G_0}(u_0)=T_0\).

Now the \(S\)-functor
\[
  P=
  \operatorname{Hom}_{S\text{-}\mathrm{gr}}
  \bigl((\mathbb Z/n\mathbb Z)_S^{\,r},G\bigr)
\]
is represented by an \(S\)-scheme of finite type: it is a closed
subscheme of \(r\) copies of
\[
  {}_nG
  =
  \operatorname{Hom}_{S\text{-}\mathrm{gr}}
  \bigl((\mathbb Z/n\mathbb Z)_S,G\bigr)
  =
  \operatorname{Ker}\bigl(G\xrightarrow{\,n\,}G\bigr).
\]
But \(P\) is smooth over \(S\) \textup{(\SGARefLink{sga3:ix:theorem:3:6}{Exposé IX, 3.6})}.
Therefore \(u_0\in P(k)\) lifts to a section \(u\in P(S)\), by
Hensel's lemma \textup{(\SGARefLink{sga3:xi:corollary:1:11}{Exposé XI, 1.11})}:
\[
  u\colon(\mathbb Z/n\mathbb Z)_S^{\,r}\longrightarrow G.
\]

Consider \(H=\operatorname{Centr}_G(u)\).  It is a closed subgroup
scheme of \(G\), by \SGARefLink{sga3:viii:examples:6:5:item:e:01}{Exposé VIII, 6.5(e)}, and its special fiber is
\(H_0=T_0\).  Moreover, \(H\) is smooth over \(S\).  Indeed, let
\(S'=\operatorname{Spec}(A')\) be an affine \(S\)-scheme, let
\[
  u'\colon({}_n\mathbb G_{m,S'})^r\longrightarrow G_{S'}
\]
be obtained from \(u\) by base change, put
\(S'_J=\operatorname{Spec}(A'/J)\), where \(J\) is a square-zero
ideal, and let \(x\in H(S'_J)\).  Since \(G\) is smooth, \(x\) lifts
to an element \(g\in G(S')\).  Then
\[
  v=\operatorname{int}(g)(u')
\]
satisfies \(v_J=u'_J\).  By \SGARefLink{sga3:ix:theorem:3:2}{Exposé IX, 3.2}, there is therefore an
element \(g'\in G(S')\) such that \(g'_J=e\) and
\(\operatorname{int}(g')(v)=u'\).  It follows that \(h=g'g\) belongs
to \(H(S')\) and satisfies \(h_J=x\).
\footnote{Editor's note: the reference to Exposé VIII, 6.5(e) was
added above, and the preceding smoothness argument has been
expanded.}

Let \(H^0\) be the identity component of \(H\).  It is a subgroup
scheme of \(G\), smooth and with connected fibers, whose special
fiber is a torus.  By \SGARefLink{sga3:x:theorem:8-1}{Exposé X, 8.1}, it is a torus, necessarily split
by \SGARefLink{sga3:x:corollary:4-6}{Exposé X, 4.6}.  Put \(T=H^0\), and let
\[
  C=\operatorname{Centr}_G(T).
\]
This is a closed subgroup of \(G\) \textup{(\SGARefLink{sga3:viii:examples:6:5:item:e:01}{Exposé VIII, 6.5(e)})}
and is smooth \textup{(\SGARefLink{sga3:xi:corollary:2:4}{Exposé XI, 2.4})}.  Consider \(C^0\)
\textup{(in fact \(C^0=C\), but we do not need this)}.  Then
\(C^0\supset T\), and both are smooth groups with connected fibers.
They coincide at \(s_0\), hence in a neighborhood of \(s_0\).
After shrinking \(S\), we may therefore suppose that \(C^0=T\);
then, a fortiori, \(T\) is maximal.

\par\medskip\noindent
\SGAInlineRefTarget{sga3:xix:remark:6-2}\textbf{\textit{\SGARefLink{sga3:xix:remark:6:2}{Remark 6.2}.}}---
\textit{The proof shows in particular that the reductive rank of
\(G_{\bar s}\) is constant in a neighborhood of \(s=s_0\).}
\SGARefTarget{sga3:xix:remark:6:2}

\par\medskip\noindent
\SGAInlineRefTarget{sga3:xix:proposition:6-3}\textbf{\textit{\SGARefLink{sga3:xix:proposition:6:3}{Proposition 6.3}.}}---
\textit{Let \(S\) be a scheme, let \(G\) be an \(S\)-group scheme
that is smooth and of finite presentation over \(S\), and let \(Q\)
be a subtorus of \(G\).}
\SGARefTarget{sga3:xix:proposition:6:3}
